\section{The Tessellation Catalog}

This appendix gives the full enumerated catalog of regular tessellations with their measured properties. Each was built and swept at a fixed scale by the engine, then read for the one structural fact that decides everything, whether the substrate carries the spinor sites. The catalog is the evidence behind the selection. It runs the whole family, dimension by dimension, and singles out exactly one mesh, the order-4 icositetrachoric honeycomb $\{3,4,3,4\}$, as the place the physics lives.

Two pieces of vocabulary are needed up front. A \emph{mesh} is the whole discrete space, one regular tessellation grown shell by shell. A \emph{dock} is a single cell of that mesh, a here-now, and a \emph{site} is one of the directions out of a dock. The candidate substrate is the mesh whose docks each carry exactly $24$ sites, and those $24$ sites are the directions that matter. When the catalog says a tessellation carries \emph{spinor sites}, it means its $24$-site dock is faceted as the $24$-cell, equivalently the $D_4$ root system, and therefore transforms with half-integer-spin pieces under its own symmetry. That is the property nothing else in the catalog can fake.

\subsection{The columns}

Each row of the catalog records a small fixed set of facts about one tessellation. The columns are read as follows.

\begin{table}[ht]
\centering
\begin{tabular}{@{}ll@{}}
\toprule
column & meaning \\
\midrule
symbol & the Schl\"afli symbol of the tessellation \\
class & compact, paracompact, or noncompact \\
crystallographic & yes if every symbol entry is $3$, $4$, or $6$ \\
spinor sites & yes if the dock is the $24$-cell, the $D_4$ root sites \\
growth & the shell growth ratio, how fast each shell outgrows the last \\
buildable & whether the engine grows it cleanly \\
\bottomrule
\end{tabular}
\caption{The catalog columns and what each records.}
\label{tab:catalog-columns}
\end{table}

A few words sharpen these. The \emph{class} is compact when the docks are bounded and the mesh has no ideal element, paracompact when the mesh has one ideal Euclidean element (a cusp), and noncompact when the docks themselves are unbounded hyperbolic pieces. The flat cusp of a paracompact mesh is the seat of physical space, so a paracompact mesh in four dimensions gives a three-dimensional flat cusp, ordinary room. A tessellation is \emph{crystallographic} only if every entry of its symbol is $3$, $4$, or $6$, the values for which a root system exists. A $5$, a $7$, or an $8$ anywhere in the symbol kills it. The \emph{growth} is the strength of the surrounding bath, how fast each new shell of docks outnumbers the one before, the openness that lets radiation fall behind the wake and never return.

The holographic correlator is the universal clean $1/r^2$ return on every buildable substrate, so it is not repeated per row. The same engine that grows every row also measures it, and it comes out the same shape everywhere.

\subsection{The headline}

\begin{result}[The spinor sites are rare and dimension-gated]
The spinor sites appear on exactly seven tessellations, all four-dimensional or five-dimensional, every one built around the $24$-cell. In two and three dimensions no tessellation carries them. Among the four-dimensional ones, $\{3,4,3,4\}$ is the unique paracompact regular with both the spinor sites and a flat three-dimensional cusp. That makes it the substrate.
\end{result}

The rest of this appendix is the table-by-table proof of that statement. The catalog is walked from two dimensions up, and the property either appears or it does not, with no room left for a borderline case.

\subsection{Two dimensions, the drawable faces}

When the bulk is two-dimensional the flat layer is a one-dimensional horocycle, so physical space is one-dimensional. This is the most degenerate case, far too thin for matter, but it is the cheapest to draw and the easiest to test holography on. All ten buildable two-dimensional regulars are compact. None carries spinor sites. The crystallographic ones are exactly those whose entries are all $3$, $4$, or $6$.

\begin{table}[ht]
\centering
\begin{tabular}{@{}llllll@{}}
\toprule
symbol & class & crystallographic & spinor sites & growth & note \\
\midrule
$\{7,3\}$ & compact & no ($7$) & no & $1.174$ & the cover image, cheapest holography testbed \\
$\{3,7\}$ & compact & no ($7$) & no & $1.174$ & dual of $\{7,3\}$ \\
$\{8,3\}$ & compact & no ($8$) & no & $1.234$ & more curved than $\{7,3\}$ \\
$\{3,8\}$ & compact & no ($8$) & no & $1.234$ & dual of $\{8,3\}$ \\
$\{5,4\}$ & compact & no ($5$) & no & $1.280$ & non-crystallographic, $5$-fold \\
$\{4,5\}$ & compact & no ($5$) & no & $1.280$ & dual of $\{5,4\}$ \\
$\{6,4\}$ & compact & \textbf{yes} & no & $1.360$ & crystallographic, $4$- and $6$-fold \\
$\{4,6\}$ & compact & \textbf{yes} & no & $1.360$ & crystallographic \\
$\{5,5\}$ & compact & no ($5$) & no & $1.400$ & self-dual, non-crystallographic \\
$\{6,6\}$ & compact & \textbf{yes} & no & $1.502$ & self-dual, crystallographic \\
\bottomrule
\end{tabular}
\caption{The two-dimensional regular tessellations, all compact, none carrying spinor sites.}
\label{tab:catalog-2d}
\end{table}

Two further two-dimensional regulars exist but the engine does not grow them, a non-convex star tiling $\{7/2,7\}$ and a paracompact apeirogonal tiling $\{\infty,3\}$ whose docks are ideal apeirogons. They are listed for completeness and left out of the buildable count.

The reading is worth stating, because two dimensions is the one place where two desirable properties coincide. Here you can have compact and crystallographic together, as $\{6,4\}$, $\{4,6\}$, and $\{6,6\}$ do. In every higher dimension those two pull apart. But even these three carry no spinor sites, and the physical space they give is only one-dimensional, far too degenerate to hold matter. $\{7,3\}$ survives the appendix not as a physics candidate but as the clean computational anchor, the cheapest mesh on which to confirm that the holographic $1/r^2$ return is really there.

\subsection{Three dimensions, the four compact regulars}

When the bulk is three-dimensional the flat layer is a two-dimensional horosphere, so physical space is two-dimensional, still under-dimensional for the world we live in. There are four compact regular three-dimensional tessellations. Every one of them is bounded and cusp-free, and every one carries a $5$ somewhere in its symbol. So all four are non-crystallographic. No root system, no gauge structure, and no spinor sites.

\begin{table}[ht]
\centering
\begin{tabular}{@{}lllllll@{}}
\toprule
symbol & class & crystallographic & spinor sites & degree & growth & note \\
\midrule
$\{5,3,4\}$ & compact & no ($5$) & no & $12$ & $1.393$ & dodecahedral, the comparison substrate \\
$\{4,3,5\}$ & compact & no ($5$) & no & $6$ & $1.393$ & dual of $\{5,3,4\}$ \\
$\{3,5,3\}$ & compact & no ($5$) & no & $20$ & $1.386$ & icosahedral, self-dual \\
$\{5,3,5\}$ & compact & no ($5$) & no & $12$ & $1.520$ & self-dual, most curved 3D compact \\
\bottomrule
\end{tabular}
\caption{The four compact three-dimensional regulars, all non-crystallographic, none carrying spinor sites.}
\label{tab:catalog-3d-compact}
\end{table}

The dodecahedral mesh $\{5,3,4\}$ is kept through the paper as the comparison substrate. Its dock has $12$ sites, and those $12$ directions decompose under icosahedral symmetry as $1 + 3 + 3' + 5$, all integer spin, no spinor anywhere. It is the cleanest, best-addressed three-dimensional mesh, the natural foil against which the spinor property is tested, but it is a foil, not a winner.

\subsection{Three dimensions, the eleven paracompact regulars}

The paracompact three-dimensional regulars each carry one ideal Euclidean element, a cusp. There are eleven. Nine are crystallographic, with entries drawn only from $3$, $4$, and $6$. Two carry a $5$ and so are not.

\begin{table}[ht]
\centering
\begin{tabular}{@{}lllll@{}}
\toprule
symbol & class & crystallographic & spinor sites & growth \\
\midrule
$\{3,3,6\}$ & paracompact & \textbf{yes} & no & $1.321$ \\
$\{6,3,3\}$ & paracompact & \textbf{yes} & no & $1.321$ \\
$\{3,4,4\}$ & paracompact & \textbf{yes} & no & $1.430$ \\
$\{4,4,3\}$ & paracompact & \textbf{yes} & no & $1.430$ \\
$\{4,3,6\}$ & paracompact & \textbf{yes} & no & $1.482$ \\
$\{6,3,4\}$ & paracompact & \textbf{yes} & no & $1.482$ \\
$\{3,6,3\}$ & paracompact & \textbf{yes} & no & $1.508$ \\
$\{5,3,6\}$ & paracompact & no ($5$) & no & $1.578$ \\
$\{6,3,5\}$ & paracompact & no ($5$) & no & $1.578$ \\
$\{4,4,4\}$ & paracompact & \textbf{yes} & no & $1.631$ \\
$\{6,3,6\}$ & paracompact & \textbf{yes} & no & $1.629$ \\
\bottomrule
\end{tabular}
\caption{The eleven paracompact three-dimensional regulars, nine crystallographic, none carrying spinor sites.}
\label{tab:catalog-3d-paracompact}
\end{table}

Three noncompact three-dimensional regulars round out the dimension, $\{3,3,7\}$, $\{7,3,3\}$, and $\{3,7,3\}$. The engine grows them, but they carry hyperbolic docks and a $7$, so they are non-crystallographic and carry no spinor sites.

The reading in three dimensions is the pattern that will repeat. You can have compact but non-crystallographic, the four regulars, all bearing a $5$. Or you can have crystallographic but paracompact, the $3$/$4$/$6$ family. You never get both at once, and either way the physical space is only two-dimensional. Crystallographic-ness bundles with an ideal element here, exactly the bundle that recurs one dimension up. $\{5,3,4\}$ is retained as the cleanest representative and the comparison substrate. Three dimensions is ruled out regardless, on the dimension count alone, because two-dimensional physical space cannot hold the matter the theory needs.

\subsection{Four dimensions, the substrate}

When the bulk is four-dimensional the flat cusp is a three-dimensional horosphere, real space, ordinary room. This is the first dimension where the cusp has the right number of dimensions for matter. The five compact four-dimensional regulars all carry a $5$, so none is crystallographic. The crystallographic four-dimensional regulars are paracompact, and two of them carry the spinor sites.

\begin{table}[ht]
\centering
\begin{tabular}{@{}lllll@{}}
\toprule
symbol & class & crystallographic & spinor sites & growth \\
\midrule
$\{3,3,3,5\}$ & compact & no ($5$) & no & $1.379$ \\
$\{5,3,3,3\}$ & compact & no ($5$) & no & $1.379$ \\
$\{4,3,3,5\}$ & compact & no ($5$) & no & $1.543$ \\
$\{5,3,3,4\}$ & compact & no ($5$) & no & $1.543$ \\
$\{5,3,3,5\}$ & compact & no ($5$) & no & $1.621$ \\
$\{4,3,4,3\}$ & paracompact & \textbf{yes} & \textbf{yes} & $1.575$ \\
$\{3,4,3,4\}$ & paracompact & \textbf{yes} & \textbf{yes} & $1.575$ \\
$\{4,4,3,3\}$ & noncompact & \textbf{yes} & no & $1.613$ \\
$\{4,4,4,4\}$ & noncompact & \textbf{yes} & no & $1.950$ \\
\bottomrule
\end{tabular}
\caption{The four-dimensional regulars. Only $\{3,4,3,4\}$ and its dual pair crystallographic structure, spinor sites, and a flat three-dimensional cusp.}
\label{tab:catalog-4d}
\end{table}

\begin{theorem}[The substrate is $\{3,4,3,4\}$]
Among the four-dimensional regular tessellations, $\{3,4,3,4\}$ and its dual $\{4,3,4,3\}$ are the only ones that carry the spinor sites and a flat three-dimensional cusp at once. The compact regulars all carry a $5$ and so carry no spinor sites. The noncompact regulars $\{4,4,3,3\}$ and $\{4,4,4,4\}$ are crystallographic but their unbounded docks carry no $24$-cell, so no spinor sites. The dock of $\{3,4,3,4\}$ has $24$ sites, and those sites are the $D_4$ roots, the $24$-cell vertices, and the unit Hurwitz quaternions. The physics lives on $\{3,4,3,4\}$.
\end{theorem}

A star-regular and an apeirogonal four-dimensional regular exist as well, $\{3,3,5,5/2\}$ among them, but the engine does not grow them and they fall outside the buildable count.

\subsection{Five dimensions, the spinor pentacombs}

The five-dimensional regulars confirm the pattern rather than competing with it. The spinor sites appear wherever a $24$-cell substructure sits inside the dock, and in five dimensions five such tessellations exist. These meshes carry spin and curvature together, so a fermion both forms and propagates on them.

\begin{table}[ht]
\centering
\begin{tabular}{@{}lllll@{}}
\toprule
symbol & class & crystallographic & spinor sites & growth \\
\midrule
$\{4,3,3,4,3\}$ & paracompact & \textbf{yes} & \textbf{yes} & $1.830$ \\
$\{3,4,3,3,4\}$ & paracompact & \textbf{yes} & \textbf{yes} & $1.830$ \\
$\{3,3,4,3,3\}$ & paracompact & \textbf{yes} & \textbf{yes} & $1.672$ \\
$\{3,4,3,3,3\}$ & paracompact & \textbf{yes} & \textbf{yes} & $1.640$ \\
$\{3,3,3,4,3\}$ & paracompact & \textbf{yes} & \textbf{yes} & $1.640$ \\
\bottomrule
\end{tabular}
\caption{The five spinor pentacombs, all paracompact, all crystallographic, all carrying spinor sites through a $24$-cell substructure.}
\label{tab:catalog-5d}
\end{table}

The pentacombs show that the spinor sites are not unique to four dimensions. They show why four dimensions is the floor. The pentacombs give a four-dimensional flat cusp, one space dimension too many for the world we measure, while $\{3,4,3,4\}$ gives the three-dimensional cusp that matches it. The five-dimensional family is the upper bracket of the same gated property, present but over-dimensioned.

\subsection{The whole catalog at a glance}

The complete count, gathered into one table, shows how narrow the spinor property is. Of every regular tessellation the engine grows across all these dimensions, only seven carry the spinor sites, and only one of those pairs them with a three-dimensional flat cusp.

\begin{table}[ht]
\centering
\begin{tabular}{@{}llll@{}}
\toprule
dimension & tessellations & crystallographic & carry the spinor sites \\
\midrule
2D & $10$ buildable & $\{6,4\}$, $\{4,6\}$, $\{6,6\}$ & none \\
3D compact & $4$ & none (all carry a $5$) & none \\
3D paracompact & $11$ & $9$ of $11$ & none \\
3D noncompact & $3$ & none (carry a $7$) & none \\
4D compact & $5$ & none (all carry a $5$) & none \\
4D paracompact & $\{4,3,4,3\}$, $\{3,4,3,4\}$ & both & \textbf{both} \\
4D noncompact & $\{4,4,3,3\}$, $\{4,4,4,4\}$ & both & none \\
5D paracompact & $5$ & all & \textbf{all five} \\
\bottomrule
\end{tabular}
\caption{The whole catalog at a glance. The spinor sites sit only on the seven $24$-cell-faceted tessellations.}
\label{tab:catalog-glance}
\end{table}

\begin{principle}[The selection]
The spinor sites are rare and dimension-gated. They sit only on the seven $24$-cell-faceted tessellations, the two in four dimensions and the five in five dimensions. Among those, $\{3,4,3,4\}$ alone pairs the spinor sites with a flat three-dimensional cusp, three space dimensions for matter. That is the selection, and it is forced, not chosen.
\end{principle}

\subsection{See also}

The selection argument that reads this catalog is worked in the tessellation survey, and dimension by dimension in the dimension ladder. The comparison substrate $\{5,3,4\}$ is taken up in the spin obstruction and in the treatment of matter on $\{5,3,4\}$. The five-dimensional pentacombs are developed in the pentacomb resolution. The builders that grow every row here are described in the computational toolkit appendix.
